% Short answers to Chapter 4's exercises, lectures/L04/appendix/b_exercises.md.

\answerchapter{c4:ch}

\answer{c4:ex:1}{Metastability in your own words}
An input changing inside the setup/hold window can leave the flip-flop's cross-coupled feedback
loop balanced between its two stable states, so its output sits at a voltage that is neither a
valid \code{0} nor a valid \code{1}. It resolves on its own, but how long that takes depends on how
nearly balanced it was left, which nobody controls and which differs every time. Until it resolves,
each downstream gate reads the voltage against its own threshold, so one can see \code{0} and
another \code{1} at the same instant, and the parts of the circuit they feed act on different
values: a counter increments in one place and not another, a state machine lands in an impossible
state.

\answer{c4:ex:2}{The setup-and-hold window}
A synchronous signal is launched by an edge of the same clock, so it arrives a fixed, known delay
after that edge and can be designed to settle before the next edge's window. A pushbutton changes
whenever a person presses it, with no relation to the clock's phase, so over enough presses some
changes will fall inside the window however the logic around it is built. The risk therefore cannot
be designed away, only contained: by giving the unresolved value time to settle before anything
uses it.

\answer{c4:ex:3}{Drawing a double-flop synchronizer}
\code{FF1} samples the asynchronous input directly, so it is the stage whose input can change inside
its window and go metastable. Its output has a whole clock period to resolve before \code{FF2}
samples it, and the chance of still being unresolved shrinks exponentially with that time, so
\code{FF2} goes wrong only in the rare case where \code{FF1} stayed unresolved for the whole period;
that is why only \code{FF2}'s output is used. The residual probability is not zero, but the
resulting mean time between failures is so long (far beyond the product's life, at these speeds)
that the risk is acceptable.

\answer{c4:ex:4}{A third synchronizer stage}
\pt{a}It turns on the required mean time between failures against the failure rate this input
produces. A third stage wins when two stages do not give an MTBF comfortably beyond what is
required: a fast-changing input, a flip-flop with poor metastability characteristics, or a failure
that must essentially never happen. Two stages are enough when their MTBF already exceeds the
requirement by orders of magnitude, and then the third only costs a cycle of latency and a
flip-flop per input.
\pt{b}\begin{itemize}
  \item A frequently changing input (high \code{f_data}) lowers MTBF in proportion and pushes
    towards a third stage; a button pressed once a second pushes away from it.
  \item A family with a slow decay (large \code{tau}) or a wide window (\code{T_w}) lowers MTBF and
    pushes towards a third stage; \code{tau} sits in the exponent, so it matters most.
  \item A costly failure (safety, money) raises the MTBF you need, pushing towards a third stage; a
    demo LED does not.
\end{itemize}
\pt{c}Two stages: a button changing a few times a second already gives an astronomically long MTBF,
and a third stage buys nothing but latency. A third stage is worth arguing for on, say, an
asynchronous line toggling at around 10\,MHz in a braking controller, where \code{f_data} and the
cost of a failure do the work.

\answer{c4:ex:5}{Assert async, release sync}
\ptnote{Assertion}Both stages go low the instant \code{reset_n} does, without waiting for an edge.
That matters when the clock is stopped (the edge never comes), unstable (its edges cannot be
trusted, as at power-up), or when the system must reach its safe state at once.
\ptnote{Recovery and removal}Releasing an asynchronous reset is itself an asynchronous change: if
the release lands inside the window around a clock edge, the flip-flop may leave reset on that edge
or the next, or go metastable.
\ptnote{Synchronous release}The design sees the release only through the two-flop chain, so
\code{reset_s2_n} rises just after a clock edge, a full period clear of the next; every flip-flop
sees the same clean release and leaves reset on the same edge, so no two of them can disagree for a
cycle about the state.

\answer{c4:ex:6}{\code{led_toggle_sync} by hand}
\pt{e}\begin{itemize}
  \item No. The first edge only loads the press into \code{button_s1_n}.
  \item Two edges: on the second, \code{button_s2_n} goes to \code{0} while \code{button_s3_n} still
    holds \code{1}, so \code{button_edge_s2} is \code{1} for one cycle. The LED toggles on the third
    edge, the same edge that ends the pulse.
  \item Yes: a small, fixed latency, two edges through the synchronizer and one more for the toggle
    flip-flop, as \secref{c4:a3} and \secref{c4:a9} predict.
\end{itemize}

\answer{c4:ex:7}{Synchronized is not debounced}
\pt{a}Synchronization and edge detection; not debouncing.
\pt{b}No more than one sampling edge fell inside the bounce, so the samples went from released to
pressed exactly once and the chain saw one clean transition. The design controls only its sampling
(the clock period), not the bounce, which the switch decides; with a 1\,s period a 1\,ms bounce
almost always fits between two edges, but nothing in the circuit makes it.
\pt{c}\begin{itemize}
  \item At most once, and usually not at all: about one press in a thousand has an edge inside the
    bounce.
  \item 1\,ms\,/\,20\,ns $= 50\,000$ times.
  \item Every sampled bounce is a fresh falling edge, so the LED toggles several times for one press
    and ends on or off unpredictably. In CircuitVerse the whole bounce fell between two edges a
    second apart and was never seen.
\end{itemize}
\pt{d}Hardware: an RC low-pass filter smooths the bounce into one slow edge, and a Schmitt-trigger
input turns that into a single clean transition. Digital: after the synchronizer, accept a new level
only once it has held for a few milliseconds' worth of cycles, restarting the count on any change;
or, with Chapter~\ref{c7:ch}'s timer, ignore further edges for a fixed interval after a press.
\pt{e}Debouncing removes the extra transitions but not the asynchrony: the debounced signal still
changes with no relation to the clock and can still land inside a flip-flop's setup/hold window, so
metastability still needs the double-flop chain.

\answer{c4:ex:8}{The reset synchronizer}
\pt{b}The reset branch forces the in-reset value, \code{'0'}, at once, while the clocked branch
shifts in the released value, \code{'1'}, which is how the release reaches the output: one branch
asserts the reset and the other releases it.
\pt{d}\begin{itemize}
  \item \code{s2} says the signal is the output of the second flip-flop of a two-flop synchronizer;
    the name would become wrong if the number of stages changed, say to three.
  \item Some module has to take the raw asynchronous reset, because a synchronized one does not
    exist until it has made it. Exactly one does, so that every flip-flop in the design is released
    by the same \code{reset_s2_n} on the same edge; two reset synchronizers could release on
    different edges.
\end{itemize}

\answer{c4:ex:9}{The button synchronizer}
\pt{b}\begin{itemize}
  \item No. From all-zeros, the released level moves through as \code{s1 s2 s3} $=$ \code{000},
    \code{100}, \code{110}, \code{111}, and \code{(not s2) and s3} is \code{0} at every step: the
    chain sees a release, and a release does not pulse.
  \item All-ones: one pulse, two edges after the reset releases. All-zeros: no pulse at all, because
    the chain already reads ``pressed'' and never sees a falling edge. All-ones is the one you want:
    the press is real, and all-zeros silently loses it.
\end{itemize}
\pt{d}\begin{itemize}
  \item An out-of-range width, above all \code{COUNT = 0} with its null-range ports, is rejected
    when the design is elaborated instead of being built into something meaningless, and the
    interface states the widths the module is meant for.
  \item Chapter~\ref{c8:ch}'s \code{seq_detect_mealy}, which has no button but still needs its
    reset synchronized. Merged, it would have had to instantiate the whole module and disable the
    button half it does not use.
\end{itemize}

\answer{c4:ex:10}{Two synchronized LED toggles}
\pt{c}Two synchronizer flip-flops bring the press to \code{button_s2_n} on the second edge, where
the edge pulse goes high, and the LED flip-flop takes it on the third. An extra pipeline stage would
move the toggle to the fourth edge, and the testbench, which expects the toggle on the third, fails
it.
\pt{d}\begin{itemize}
  \item It makes each press safe from metastability, but does nothing about bounce: at 50\,MHz every
    bounce is still its own pulse, so one press can still toggle the LED several times.
  \item A physical button is asynchronous to the 50\,MHz clock, and the raw one reached two
    flip-flops on the same edge, \code{button_prev_n} directly and \code{led_s} through the edge
    logic; a press inside the setup/hold window could leave either metastable, or be read as pressed
    by one and released by the other, losing the toggle or doubling it.
  \item Metastability: a simulated flip-flop always captures a clean \code{'0'} or \code{'1'}.
    Bounce, by contrast, a testbench could at least imitate by driving a burst of transitions.
\end{itemize}
